\chapter{物理动机与相关理论背景}

\section{重夸克偶素联合产生的研究意义}

重夸克偶素末态兼具微扰和非微扰强子化效应，是检验夸克偶素产生机制的重要对象。
对于多粒子联合产生过程，仅依靠单部分子散射往往难以完整解释其产额和运动学分布，
因此多部分子相互作用的贡献需要被系统考虑。

\(J/\psi + J/\psi + \phi\) 末态一方面包含双 \(J/\psi\) 子系统，另一方面又引入了奇异夸克相关的 \(\phi\) 介子，
使其在末态组合、背景组成和理论解释上都区别于单纯的双 \(J/\psi\) 或三 \(J/\psi\) 分析。
这也是本文将物理动机独立成章，而不直接并入绪论的原因。

\section{与公开 CMS 分析结构的对应关系}

\texttt{struct-ref/} 中的公开 CMS quarkonium 相关文章显示，成熟分析通常会把“理论动机”和
“分析对象的测量定义”放在实验细节之前，以便读者先明确：

\begin{itemize}
  \item 目标过程对应的物理问题是什么；
  \item 需要报告的观测量是什么；
  \item 结果是积分结果、微分结果，还是显著性与上限结果。
\end{itemize}

对本文而言，后续需要在这一章内补全的内容主要包括：

\begin{itemize}
  \item \(J/\psi + J/\psi + \phi\) 过程与多部分子相互作用研究的关系；
  \item 可用的理论产生机制或现象学解释框架；
  \item 本分析最终采用的 fiducial 定义与观测量定义。
\end{itemize}

\section{本文采用的分析叙事框架}

为了和后续实验章节保持一致，本文采用如下分析叙事框架：
先定义物理问题与观测对象，再给出数据样本与候选选择，随后说明效率和拟合方法，
最后报告结果并讨论其与相关工作的联系。

本章暂不写入任何未经最终确认的理论数值、预测截面或比较结论。
若后续需要加入公开理论结果，应在确认来源和适用条件后再写入正文，并同步补充参考文献。
